\chapter{Mise en œuvre I : le rendu d'images}

Première application complète, pas à pas. C'est ici que Monte-Carlo n'est pas
une astuce mais une nécessité absolue.

\section{Le problème : une équation qu'on ne peut pas résoudre}

La lumière qui repart d'un point $x$ vers l'œil, dans la direction
$\omega_o$, vaut ce que la surface émet elle-même plus tout ce qu'elle reçoit
et renvoie :

\[ L_o(x,\omega_o) \;=\; L_e(x,\omega_o) \;+\;
   \int_{\Omega} L_i(x,\omega_i)\; f_r(x,\omega_i,\omega_o)\; \cos\theta_i \; d\omega_i \]

En français, terme par terme :

\begin{itemize}
  \item $L_e$ : ce que la surface \textbf{émet} (nul, sauf si c'est une lampe) ;
  \item $\int_{\Omega}$ : la somme sur \textbf{toutes les directions} de
        l'hémisphère au-dessus du point — une infinité de directions ;
  \item $L_i(x,\omega_i)$ : la lumière qui \textbf{arrive} de la direction
        $\omega_i$ ;
  \item $f_r$ : la \textbf{BRDF}, qui dit quelle fraction de ce qui arrive
        d'une direction repart vers une autre — la matière ;
  \item $\cos\theta_i$ : la lumière rasante étale son énergie sur une plus
        grande surface, donc elle compte moins. C'est le même effet qui fait
        qu'il fait plus froid en hiver.
\end{itemize}

Deux raisons rendent cette équation insoluble : elle est \textbf{récursive}
($L_i$ est le $L_o$ d'un autre point, qui dépend d'un autre encore), et son
domaine est \textbf{continu et infini}.

\section{La transformation en algorithme}

Appliquons la formule du chapitre 2 : au lieu de sommer sur toutes les
directions, tirons-en quelques-unes au hasard.

\begin{nkretenir}[Les six pas, dans l'ordre]
\begin{enumerate}
  \item \textbf{Tirer} une direction $\omega_i$ selon une densité $p$ choisie
        par nous ;
  \item \textbf{lancer} un rayon dans cette direction et trouver ce qu'il
        touche ;
  \item \textbf{récupérer} la lumière qui vient de là — par récursion ;
  \item \textbf{pondérer} : multiplier par $f_r \cdot \cos\theta_i$, puis
        \textbf{diviser par $p(\omega_i)$} ;
  \item \textbf{accumuler} dans la somme du pixel ;
  \item \textbf{recommencer} $N$ fois, puis diviser par $N$.
\end{enumerate}
\end{nkretenir}

\section{Le choix de $p$, et la simplification qui en découle}

Tirer uniformément sur l'hémisphère fonctionne, mais bruite : on gaspille des
rayons dans les directions rasantes, que le $\cos\theta_i$ va presque annuler.

L'échantillonnage \textbf{en cosinus} tire donc plus souvent près de la normale,
avec la densité
\[ p(\omega_i) = \frac{\cos\theta_i}{\pi} \]

Regardez ce qui se passe quand on remplace, pour une surface diffuse dont la
BRDF vaut $f_r = \rho/\pi$ :

\[ \frac{f_r \cdot \cos\theta_i}{p(\omega_i)}
 = \frac{\dfrac{\rho}{\pi}\cos\theta_i}{\dfrac{\cos\theta_i}{\pi}}
 = \rho \]

\begin{nkretenir}[Pourquoi tous les moteurs font ça]
Le $\cos\theta_i$ \textbf{se simplifie} avec celui de la densité, et le $\pi$
aussi. Il ne reste que l'albédo. On a donc, en même temps :
\begin{itemize}
  \item moins de bruit (on tire là où ça compte) ;
  \item moins de calcul (deux termes disparaissent algébriquement).
\end{itemize}
C'est le prototype du bon échantillonnage d'importance : celui qui rend le
code \emph{plus simple} et non plus compliqué.
\end{nkretenir}

\begin{nkcode}{Tirage en cosinus sur l'hemisphere (methode de Malley)}
// On tire un point uniformement dans le DISQUE unite, puis on le releve sur
// l'hemisphere. C'est exactement equivalent a un tirage en cosinus, et ca ne
// coute qu'une racine carree -- resultat classique, et bien plus rapide que
// d'inverser la repartition en angles spheriques.
NkVec3 EchantillonnerCosinus(const NkVec3& normale, float32* pdfOut) {
    const float32 u1 = math::NkRand.NextFloat32();
    const float32 u2 = math::NkRand.NextFloat32();

    const float32 r = NkMath::Sqrt(u1);
    const float32 phi = 2.f * NkMath::kPi * u2;

    const float32 x = r * NkMath::Cos(phi);
    const float32 y = r * NkMath::Sin(phi);
    const float32 z = NkMath::Sqrt(NkMath::Max(0.f, 1.f - u1));  // = cos(theta)

    // Base orthonormee autour de la normale (methode de Duff : sans branche
    // conditionnelle, et numeriquement stable meme quand n.z vaut -1).
    const float32 sign = (normale.z >= 0.f) ? 1.f : -1.f;
    const float32 a = -1.f / (sign + normale.z);
    const float32 b = normale.x * normale.y * a;
    const NkVec3 t(1.f + sign * normale.x * normale.x * a, sign * b, -sign * normale.x);
    const NkVec3 s(b, sign + normale.y * normale.y * a, -normale.y);

    if (pdfOut)
        *pdfOut = z / NkMath::kPi;          // p = cos(theta) / pi

    return t * x + s * y + normale * z;
}
\end{nkcode}

\section{Le chemin complet}

\begin{nkcode}{Contribution d'UN chemin -- la moyenne de N appels donne le pixel}
NkVec3 TracerChemin(const NkRay& rayon, nk_uint32 profondeur) {
    NkHit hit;
    if (!mScene.Intersect(rayon, hit))
        return CielVersLumiere(rayon.direction);   // le rayon part a l'infini

    NkVec3 radiance = hit.material.emission;       // la surface brille-t-elle ?

    // ROULETTE RUSSE. Couper net a une profondeur fixe perdrait toute la
    // lumiere des rebonds suivants : l'image s'assombrirait de facon
    // systematique (un BIAIS, pas du bruit). Ici on continue avec la
    // probabilite q et on divise par q : l'esperance reste exacte.
    float32 q = 1.f;
    if (profondeur > 3) {
        q = 0.8f;
        if (math::NkRand.NextFloat32() > q)
            return radiance;                       // le chemin meurt ici
    }

    // (1) tirer une direction, et connaitre SA densite
    float32 pdf = 0.f;
    const NkVec3 wi = EchantillonnerCosinus(hit.normal, &pdf);
    if (pdf <= 0.f)
        return radiance;                           // garde : jamais diviser par 0

    // (2) et (3) : suivre le rayon, recuperer ce qui vient de la
    const NkVec3 entrant = TracerChemin(NkRay(hit.position, wi), profondeur + 1);

    // (4) ponderer : BRDF * cos, DIVISE par la densite et par q.
    // Avec le tirage en cosinus, cos et pi se simplifient : il ne reste que
    // l'albedo. On ecrit quand meme la forme complete -- elle reste juste si
    // l'on change de densite un jour, et c'est la que se cachent les bugs.
    const float32 cosTheta = NkMath::Max(0.f, NkVec3::Dot(hit.normal, wi));
    const NkVec3 brdf = hit.material.albedo * (1.f / NkMath::kPi);

    radiance += entrant * brdf * cosTheta / (pdf * q);
    return radiance;
}

// (5) et (6) : la moyenne sur le pixel.
NkVec3 RendrePixel(nk_uint32 px, nk_uint32 py, nk_uint32 echantillons) {
    NkVec3 somme(0.f, 0.f, 0.f);
    for (nk_uint32 s = 0; s < echantillons; ++s) {
        // Jitter dans le pixel : sans lui, tous les rayons passent par le
        // meme point et les bords sont crenele. Avec, l'anticrenelage est
        // GRATUIT -- il vient du meme hasard qui sert a l'eclairage.
        const float32 jx = math::NkRand.NextFloat32();
        const float32 jy = math::NkRand.NextFloat32();
        somme += TracerChemin(mCamera.RayonVers((float32)px + jx, (float32)py + jy), 0);
    }
    return somme / (float32)echantillons;
}
\end{nkcode}

\begin{nkpiege}[Trois erreurs qui donnent une image plausible mais fausse]
\begin{enumerate}
  \item \textbf{oublier de diviser par la densité} : l'image reste jolie, mais
        l'éclairage est faux. Rien ne le signale — c'est le bug nº1 du path
        tracing ;
  \item \textbf{couper les chemins à profondeur fixe} : image trop sombre,
        d'autant plus que la scène a d'inter-réflexions ;
  \item \textbf{oublier le jitter dans le pixel} : bords crénelés, qu'aucun
        nombre d'échantillons ne corrigera jamais — tous les rayons passent
        exactement au même endroit.
\end{enumerate}
\end{nkpiege}

\section{Pourquoi le rendu direct des lumières change tout}

Avec le code ci-dessus, une petite lampe intense est un cauchemar : les rayons
tirés au hasard la manquent presque toujours, et les rares qui la touchent
rapportent une énergie énorme. Résultat : des \textbf{pixels blancs isolés} sur
une image sombre, ce qu'on appelle des \emph{fireflies}.

La parade s'appelle \emph{next event estimation} : à chaque rebond, on tire
\textbf{en plus} un point directement sur une source de lumière, on vérifie
qu'il est visible, et on ajoute sa contribution. On combine ensuite les deux
estimations (celle par la BRDF et celle par la lumière) avec des poids —
c'est l'\emph{échantillonnage à importances multiples}, la brique suivante.

\begin{nkexo}[Voir le bruit diminuer]
Rendez une scène simple (un plan, une sphère, un ciel uniforme) à 1, 4, 16,
64 puis 256 échantillons par pixel, en gardant les images.

Vérifiez la loi du chapitre 3 : entre 16 et 64 échantillons — quatre fois plus
— le grain doit diminuer d'un facteur 2 environ, pas 4. Puis ajoutez le tirage
direct des lumières et comparez à nombre d'échantillons égal.
\end{nkexo}
